\begin{center}
\textsc{Gauss} an \textsc{Gerling}.\quad Göttingen, 5.~September 1823.
\end{center}

\par\noindent\hspace*{\parindent}\mbox{.\ .\ .\ .\ .\ .}\quad Ich wünsche unbedingt, dass Sie den Heliotrop [auf dem
Inselsberg] genau in einer Linie mit dem Standort des Theodoliten aufstellen und das Zelt
so symmetrisch wie möglich aufschlagen. Wenn Sie der dem Brocken
zugewandten Seite eine dunkle Farbe geben könnten, wäre es umso besser. Ich
\textls[120]{hoffe}, dass der Heliotrop dort das Zelt bei Weitem überstrahlen wird. Ich
habe aber 1822 den Fall erlebt, dass das \textls[120]{weiße} Zelt in Scharnhorst meine
Beobachtungen in Garssen einen ganzen Tag lang unbrauchbar machte. Dabei
ist allerdings zu beachten, dass das Zelt mir erstens seine von der Sonne
beleuchtete Südseite zuwandte und zweitens das Heliotroplicht sehr schwach war.
Beides entfällt hier. Vorsichtshalber möchte ich den Heliotrop dennoch
\textls[120]{nicht außermittig} aufgestellt wissen, da das Zelt sonst ein wenig dazu
beitragen könnte, den Winkel zu verfälschen. Ich vermute, dass, wenn z.~B.\ das
Zelt \(\frac{1}{40}\) so hell ist wie der Heliotrop und dieser so weit entfernt steht, dass
der Abstand \(4^{\prime\prime}\) beträgt, der Winkel um \(0\mathord{\ooalign{\hfil\(\vphantom{0}^{\prime\prime}\)\hfil\cr\hfil\({,}\)\hfil\cr}}1\) verfälscht wird.\\*
\mbox{.\ .\ .\ .\ .\ .}\par
